% Figure: three line-load distributions (uniform, triangular, trapezoidal)
% with their equivalent resultant point loads located at the area centroid.
\begin{figure}[htbp]
\centering
\begin{tikzpicture}[>=latex, x=1.0cm, y=1.0cm]
  % ==== uniform ====
  \begin{scope}
    \draw[thick] (0,0) -- (3,0);
    \fill[engblue!18, draw=engblue, thick] (0,0) rectangle (3,1.0);
    % load arrows
    \foreach \x in {0.3,0.8,...,2.9} { \draw[->, engblue, thin] (\x,1.0) -- (\x,0.1); }
    % resultant at midspan
    \draw[->, very thick, engred] (1.5,1.9) -- (1.5,0.1);
    \node[engred, anchor=south, font=\scriptsize] at (1.5,1.9) {$W=w_0L$};
    \node[font=\scriptsize, anchor=north] at (1.5,-0.1) {$\bar x = L/2$};
    \node[font=\scriptsize] at (1.5,-0.7) {uniform};
    \node[engblue, font=\scriptsize, anchor=west] at (0.05,0.55) {$w_0$};
  \end{scope}
  % ==== triangular ====
  \begin{scope}[xshift=4.5cm]
    \draw[thick] (0,0) -- (3,0);
    \fill[englime!22, draw=englime!70!black, thick] (0,0) -- (3,0) -- (3,1.3) -- cycle;
    \foreach \x in {0.5,1.0,...,2.9} {
      \pgfmathsetmacro\h{\x/3*1.3}
      \draw[->, englime!60!black, thin] (\x,\h) -- (\x,0.1);
    }
    \draw[->, very thick, engred] (2.0,1.9) -- (2.0,0.1);
    \node[engred, anchor=south, font=\scriptsize] at (2.0,1.9) {$W=\tfrac12 w_{\max}L$};
    \node[font=\scriptsize, anchor=north] at (2.0,-0.1) {$\bar x = 2L/3$};
    \node[font=\scriptsize] at (1.5,-0.7) {triangular};
    \node[font=\scriptsize, anchor=west] at (3.05,1.25) {$w_{\max}$};
  \end{scope}
  % ==== trapezoidal ====
  \begin{scope}[xshift=9.0cm]
    \draw[thick] (0,0) -- (3,0);
    \fill[engamber!22, draw=engamber, thick] (0,0) -- (3,0) -- (3,1.3) -- (0,0.5) -- cycle;
    \foreach \x in {0.3,0.8,...,2.9} {
      \pgfmathsetmacro\h{0.5+\x/3*0.8}
      \draw[->, engamber!80!black, thin] (\x,\h) -- (\x,0.1);
    }
    \draw[->, very thick, engred] (1.75,1.9) -- (1.75,0.1);
    \node[engred, anchor=south, font=\scriptsize] at (1.75,1.9) {$W=\bar w L$};
    \node[font=\scriptsize, anchor=north] at (1.75,-0.1) {$\bar x$};
    \node[font=\scriptsize] at (1.5,-0.7) {trapezoidal};
    \node[font=\scriptsize, anchor=east] at (-0.05,0.5) {$w_1$};
    \node[font=\scriptsize, anchor=west] at (3.05,1.25) {$w_2$};
  \end{scope}
\end{tikzpicture}
\caption[Equivalent point loads for line distributions]{Three line-load
distributions and the single point load that reproduces each one's
resultant force and moment. The equivalent load magnitude is the area
under the loading diagram and its line of action passes through that
area's centroid. A uniform load $w_0$ over span $L$ collapses to
$W=w_0L$ at midspan. A triangular load rising to $w_{\max}$ collapses to
$W=\tfrac12 w_{\max}L$ at two-thirds of the span from the zero end, where
the triangle's centroid sits. A trapezoid is the sum of a rectangle and
a triangle, so its resultant and centroid follow from adding the two
constituent areas, which is the working engineer's method for any
piecewise-linear distribution: break it into rectangles and triangles,
sum the areas, and take the area-weighted average of the constituent
centroids. The equivalence holds for rigid-body equilibrium only; the
internal stress depends on where the load actually acts.}
\label{fig:vol03ch02:distributed-loads}
\end{figure}
